\section{Особенности реализации}

\subsection{Общий поток выполнения}

Проект собирается через CMake. Точка входа~-- единый исполняемый файл
с консольным меню. Все модули лабораторных зависят только от \texttt{common};
\texttt{common} ничего не знает про лабораторные.

Пользователь выбирает пункт меню, \texttt{main.cpp} вызывает нужный метод
\texttt{LabXUI}, который читает параметры, запускает алгоритм и выводит
результат в консоль.

\subsection{Модуль common}

Модуль содержит реализацию графа, генерацию случайных чисел по распределению Вейбулла и генерацию графов, и функции для вывода общих структур. Этот модуль отделен от лабораторных, для переиспользования кода в разных работах.

\subsubsection{Graph\textless T\textgreater: шаблонный граф}

\texttt{Graph<T>}~-- шаблонный класс для хранения графа; тип весов рёбер
задаётся параметром \texttt{T}.

\begin{lstlisting}[style=cstyle, caption={Объявление Graph (graph.h)}]
template <typename T>
class Graph {
  private:
    int numVertices;
    vector<vector<int>> adjMatrix;
    vector<vector<T>> weightMatrix;
  public:
    Graph(int n);
    void addEdge(int from, int to, T weight);
    bool hasEdge(int from, int to) const;
    T getEdge(int from, int to) const;
    int getNumVertices() const;
    vector<vector<int>> getAdjMatrix() const;
    vector<vector<T>> getWeightMatrix() const;
    Graph<T> reorder(const vector<int>& order) const;
    // ...
};
\end{lstlisting}

\paragraph{Хранение данных.} Граф хранится в двух матрицах размера $n \times n$.
\texttt{adjMatrix[i][j] = 1} означает, что ребро $(i, j)$ есть;
\texttt{weightMatrix[i][j]} хранит его вес.
Проверка наличия ребра и доступ к весу работают за $O(1)$, но памяти
уходит $O(n^2)$ независимо от числа рёбер.

Также реализованны базовые методы для добавления рёбер, получения количества вершин и доступа к матрицам.

\subsubsection{WeibullDistribution: генерация случайных чисел}

Класс хранит параметры распределения Вейбулла: масштаб $a$ и форму $c$;
по умолчанию $a = 1, c = 0.2$.

\begin{lstlisting}[style=cstyle, caption={Генерация числа по распределению Вейбулла (weibull.cpp)}]
double WeibullDistribution::generate() {
    double u = distribution(generator);
    double raw = a * std::pow(-std::log(1 - u), 1.0 / c);
    return (raw > 0) ? raw : -raw;
}
\end{lstlisting}

Метод берёт равномерное $u \in (0,1)$ из \texttt{std::mt19937}
и считает $X = a(-\ln(1 - u))^{1/c}$. Модуль нужен, чтобы результат
всегда был положительным.

\subsubsection{Generator\textless T\textgreater: генерация графов}

\texttt{Generator<T>} держит два объекта \texttt{WeibullDistribution}:
один для структуры графа (\texttt{distributionForStructure}), \\другой
для весов (\texttt{distributionForWeights}).

\begin{lstlisting}[style=cstyle, caption={Генератор весов и Вейбуллово перемешивание (generator.cpp)}]
template <typename T>
T Generator<T>::generateWeight(WeightType weightType) {
    double weight = distributionForWeights.generate();
    if (weightType == WeightType::NEGATIVE)
        return static_cast<T>(-weight);
    else if (weightType == WeightType::MIXED)
        return isNegative() ? static_cast<T>(-weight)
                            : static_cast<T>(weight);
    return static_cast<T>(weight);
}

template <typename T>
void Generator<T>::weibullShuffle(vector<int>& vertices) {
    int n = vertices.size();
    for (int i = n - 1; i > 0; i--) {
        int j = randomIndex(i + 1);
        swap(vertices[i], vertices[j]);
    }
}
\end{lstlisting}

Метод \texttt{generateGraph()} строит граф в два прохода.
В первом проходе строится остовное дерево: для вершины $i \geq 1$
предок выбирается случайно из $\{0, \ldots, i-1\}$~-- это гарантирует
связность и отсутствие циклов.
Во втором проходе для каждой вершины с оставшимся бюджетом степени
берутся все вершины $j > i$, у которых ещё нет ребра от $i$. Список
перемешивается через \texttt{weibullShuffle} и добавляются первые $d_i$ рёбер.
